We prove the proposed support-respecting statement.  All stochastic orders below are uniform because \(\mathcal D\), the treatment set, all basis dimensions, and the fold count are finite and fixed.

Fix \((d,a)\), and let \(n\) be the size of a relevant training source-arm sample.  The new sampling premise makes the observations within that source independent and identically distributed.  Bounded bases, \(E_P(|Y|^{4+\delta}\mid G=O)<\infty\), source-arm positivity, and the proportional-size condition therefore give, entry by entry,
\[
 \widehat M^h_{d,a}-M^h_{d,a}=O_P(n^{-1/2}),
 \qquad
 \widehat g^h_{d,a}-g^h_{d,a}=O_P(n^{-1/2}),                 \tag{8}
\]
by Chebyshev's inequality.  The same conclusion holds for \(\widehat M^q_{d,a}-M^q_{d,a}\) and \(\widehat g^q_{d,a}-g^q_{d,a}\).  The latter two estimates may use different sources, but mutual independence of the source samples and the union bound preserve the joint order.  A finite union bound makes (8) simultaneous over designs, arms, and folds.

On the event \(\|\widehat M^h_{d,a}-M^h_{d,a}\|_{\mathrm{op}}<c_0/2\),
\[
 \|\widehat M^h_{d,a}x\|
 \ge \|M^h_{d,a}x\|
      -\|\widehat M^h_{d,a}-M^h_{d,a}\|_{\mathrm{op}}\|x\|
 \ge \frac{c_0}{2}\|x\|.                                  \tag{9}
\]
Hence \(\widehat M^h_{d,a}\) is nonsingular and \(\|\widehat M^{h,-1}_{d,a}\|_{\mathrm{op}}\le2/c_0\); the same argument applies to \(\widehat M^q_{d,a}\).  These events hold jointly with probability tending to one.  On them, the exact population equations yield
\[
 \widehat\beta_{d,a}-\beta_{d,a}
 =\widehat M_{d,a}^{h,-1}
 \left\{(\widehat g^h_{d,a}-g^h_{d,a})
 -(\widehat M^h_{d,a}-M^h_{d,a})\beta_{d,a}\right\},       \tag{10}
\]
with the analogous identity for \(\widehat\gamma_{d,a}-\gamma_{d,a}\).  The population coefficient vectors are bounded uniformly over the finite collection: the population inverses have norm at most \(c_0^{-1}\), while the outcome and selection right-hand sides are bounded by the moment and basis envelopes.  Consequently,
\[
 \max_{d,a,\ell}\left\{
 \|\widehat\beta_{d,a}^{(-\ell)}-\beta_{d,a}\|
 +\|\widehat\gamma_{d,a}^{(-\ell)}-\gamma_{d,a}\|
 \right\}=O_P(n^{-1/2}).                                    \tag{11}
\]
The stipulated zero value on singular empirical matrices is asymptotically irrelevant.

Bounded bases turn (11) into the pointwise envelopes
\[
 |\widehat h_d-h_d|\le C\|\widehat\beta_d-\beta_d\|,
 \qquad
 |\widehat q_d-q_d|\le C\|\widehat\gamma_d-\gamma_d\|.    \tag{12}
\]
Taking \(n\asymp B\), the first envelope gives
\[
 \|\widehat h_{d,B}^{(-\ell)}-h_d\|_{2,B}^{\circ}
 =O_P(B^{-1/2}).                                               \tag{13}
\]
Conditional expectation is an \(L_2\)-contraction on the observational source, so
\[
 \|T_d(\widehat h_{d,B}^{(-\ell)}-h_d)\|_{2,O,q,d}
 \le \|\widehat h_{d,B}^{(-\ell)}-h_d\|_{2,O,h,d}
 =O_P(B^{-1/2}).                                               \tag{14}
\]
The second envelope directly gives
\[
 \|\widehat q_{d,B}^{(-\ell)}-q_d\|_{2,O,q,d}
 =O_P(B^{-1/2}).                                               \tag{15}
\]
Moreover, every source-defined version of \(T_d^\star(\widehat q_d-q_d)\) is bounded in absolute value by the right-hand envelope in (12), because conditional expectation preserves a pointwise bound.  Therefore
\[
 \|T_d^\star(\widehat q_{d,B}^{(-\ell)}-q_d)\|_{2,B}^{\circ}
 =O_P(B^{-1/2}).                                               \tag{16}
\]
Equations (13)--(16) prove (5)--(6), and either weak--strong product is \(O_P(B^{-1})=o_P(B^{-1/2})\).

Write \(\Delta h_{B,\ell}=\widehat h_{d,B}^{(-\ell)}-h_d\) and \(\Delta q_{B,\ell}=\widehat q_{d,B}^{(-\ell)}-q_d\).  The population bridges are bounded because their bases and coefficients are bounded.  Equations (11)--(12), the finite second moment of \(Y\), and Cauchy--Schwarz imply
\[
 \|\mathbf1\{G=O\}\Delta q_{B,\ell}(Y-h_d)\|_{2,B}^{\circ}
 +\|\mathbf1\{G=O\}q_d\Delta h_{B,\ell}\|_{2,B}^{\circ}
 =O_P(B^{-1/2}),                                               \tag{17}
\]
and
\[
 \|\mathbf1\{G=O\}\Delta q_{B,\ell}\Delta h_{B,\ell}\|_{2,B}^{\circ}
 =O_P(B^{-1}).                                                 \tag{18}
\]
Thus the multiplier expression is \(O_P(B^{-1/2})=o_P(1)\).

We next isolate the experimental centering coordinate omitted by the earlier argument.  For an experimental evaluation sample in arm \(a\), let \(P_{E,a}\) denote population arm-conditional expectation and \(\widehat P_{E,a}\) its empirical analogue.  For every fold-external outcome-bridge fit, the identity
\[
 \widehat\mu_{d,a}^{(-\ell)}-\mu_{d,a}
 =P_{E,a}(\widehat h_d^{(-\ell)}-h_d)
 +(\widehat P_{E,a}-P_{E,a})h_d
 +(\widehat P_{E,a}-P_{E,a})(\widehat h_d^{(-\ell)}-h_d)    \tag{19}
\]
is exact.  The first term is \(O_P(n^{-1/2})\) by (11) and bounded population basis means.  The second is \(O_P(n^{-1/2})\) by independent and identically distributed experimental sampling and boundedness of \(h_d\).  For the third term, write the fitted difference as the bounded basis multiplied by \(\widehat\beta_{d,a}^{(-\ell)}-\beta_{d,a}\).  The coefficient difference is \(O_P(n^{-1/2})\) by (11), and the empirical-minus-population basis mean is \(O_P(n^{-1/2})\).  The outcome-bridge fit is trained on observational data and is independent of the experimental evaluation sample, so their product is \(O_P(n^{-1})\).  Therefore, uniformly over the finite designs, arms, and folds,
\[
 \widehat\mu_{d,a}^{(-\ell)}-\mu_{d,a}=O_P(n^{-1/2}).        \tag{20}
\]
In particular, the load-bearing term \(P_{E,a}(\widehat h-h)\) has not been discarded.

The empirical arm probabilities obey \(\widehat p_{g,a}-p_{g,a}=O_P(n^{-1/2})\), and source-arm positivity makes their inverses stochastically bounded.  For one fold, the experimental score difference has the exact form
\[
\begin{aligned}
 \widehat\phi_{E,d}-\phi_{E,d}
 =\sum_{a=0}^1(-1)^{1-a}\mathbf1\{A=a\}\bigg[&
 \left(\widehat p_{E,a}^{-1}-p_{E,a}^{-1}\right)
 \{h_d(a)-\mu_{d,a}\}\\
 &+\widehat p_{E,a}^{-1}
 \{\Delta h_{d,a}-(\widehat\mu_{d,a}^{(-\ell)}-\mu_{d,a})\}
 \bigg],                                                     \tag{21}
\end{aligned}
\]
where the suppressed arguments agree with the theorem's score definition.  Equations (11)--(12) and (20), together with the arm-probability rate, imply
\[
 \|\widehat\phi_{E,d}-\phi_{E,d}\|_2=O_P(n^{-1/2}).         \tag{22}
\]
Thus the experimental score error is controlled jointly by coefficient error, arm-probability error, and the full three-term centering error.

For the observational score, direct subtraction gives terms proportional to the arm-probability error and to
\[
 \Delta q_d(Y-h_d)-q_d\Delta h_d-\Delta q_d\Delta h_d.       \tag{23}
\]
Boundedness of \(q_d,h_d\), (11)--(12), and \(E_P(Y^2\mid G=O)<\infty\) imply
\[
 \|\widehat\phi_{O,d}-\phi_{O,d}\|_2=O_P(n^{-1/2}).         \tag{24}
\]
The experimental population score is bounded, and the observational population score has a finite fourth moment because \(q_d,h_d\) are bounded and \(E_P(|Y|^{4+\delta}\mid G=O)<\infty\).

Let \(\mathbb P_n\) be an evaluation-source empirical measure.  For either source and any fixed design,
\[
\begin{aligned}
 |\mathbb P_n\widehat\phi_{g,d}^2-P\phi_{g,d}^2|
 \le{}&| (\mathbb P_n-P)\phi_{g,d}^2|\\
 &+\{\mathbb P_n(\widehat\phi_{g,d}-\phi_{g,d})^2\}^{1/2}
   \{\mathbb P_n(\widehat\phi_{g,d}+\phi_{g,d})^2\}^{1/2}. \tag{25}
\end{aligned}
\]
The first term is \(O_P(n^{-1/2})\) by Chebyshev's inequality applied to \(\phi_{g,d}^2\).  Equations (22)--(24), their bounded-basis envelopes, and the empirical second-moment bound for \(Y\) make the second term \(O_P(n^{-1/2})\), even though the empirical arm proportions and centering means are computed from the evaluation sample.  Fixed-fold averaging and a finite union bound preserve the rate.  With \(n\asymp m\) in the pilot and \(n\asymp B\) in the main study, (25) yields
\[
 \max_{d\in\mathcal D}\max\bigl\{
 |\widehat V_{E,d}-V_{E,d}|,
 |\widehat V_{O,d}-V_{O,d}|\bigr\}=O_P(m^{-1/2})             \tag{26}
\]
and, first uniformly in \(d\) and hence also at the pilot-selected random design,
\[
 \max_{g\in\{E,O\}}
 |\widetilde V_{g,\widehat d_m,B}-V_{g,\widehat d_m}|
 =O_P(B^{-1/2}).                                               \tag{27}
\]

Equations (13)--(18) and (26)--(27) verify all four bridge rates, both product conditions, the multiplier condition, and both variance-rate conditions with
\[
 r_{V,m}=m^{-1/2},\qquad r_{\widetilde V,B}=B^{-1/2}.
\]
They also verify the associated vanishing conditions.  The pilot is independent of the main study by the newly stated premise, and the two main-study sources have the independent and identically distributed sampling structure required by Theorem~\(\mathrm{thm{:}postselection\text{-}wald}\).  We may therefore apply Theorem~\(\mathrm{thm{:}pilot\text{-}oracle\text{-}regret}\) and Theorem~\(\mathrm{thm{:}postselection\text{-}wald}\), which give the asserted oracle regret, selected-design normal limit, and Wald coverage.  The argument depends on the fixed-dimensional basis, rank, moment, positivity, sampling, and fold-size conditions stated in the proposed theorem, so it does not extend these conclusions to all of \(\mathfrak P_{\mathcal D}^{\mathrm{obs}}\).